\subsection{\texorpdfstring{第~\ref{sec:topological-manifolds}~节}{拓扑流形}}\label{subsec:solutions-topological-manifolds}

\begin{description}[font=\normalfont,leftmargin=2em,labelwidth=1.5em]
\RGsolution{ex:recognize-topological-manifolds}{recognize-topological-manifolds}
在下面各例中, 拓扑流形定义的条件 (i) 和 (iii) 都从周围空间继承而来, 所以只需考虑条件 (ii).
\begin{enumerate}[label=(\alph*),ref=\alph*,leftmargin=*]
\item $D^2$ 是 $\mathbb{R}^2$ 的开子集, 因而显然是拓扑 2-流形.
\item $M=S^2\setminus\{p\}$ 是 2 维拓扑流形. 把 $\mathbb{R}^2$ 看成复平面, 并不妨设 $p$ 是球面的北极, 则定义
\[
  \varphi(z)=\left(\frac{2x}{1+|z|^2},\frac{2y}{1+|z|^2},\frac{|z|^2-1}{1+|z|^2}\right),
\]
其中 $z=x+iy$, 就得到一个简单的同胚 $\varphi:\mathbb{R}^2\cong\mathbb{C}\to M$. 这个映射的逆称为球极投影 [参见 Exercise~\ref{ex:exercise-differentiable-structures}(\ref{ex:stereographic-sphere-atlas})].
\item $N=S^2\setminus\{p,q\}$ 也是拓扑 2-流形. 不妨设 $p$ 和 $q$ 分别为球面的北极和南极, 则 (b) 中定义的 $\varphi$ 是开集 $\mathbb{R}^2\setminus\{0\}\cong\mathbb{C}\setminus\{0\}$ 与 $N$ 之间的同胚.
\item 圆柱面 $V=\{(x,y)\in\mathbb{R}^3\mid x^2+y^2=1\}$ 是拓扑 2-流形. 对每个点 $p:=(x,y,z)\in V$, 取法向量 $n:=(x,y,0)$, 并考虑由两条坐标轴生成且不与 $n$ 平行的平面. 则 $V$ 中 $p$ 的某个邻域同胚于该邻域在此平面上的投影.
\item 锥面 $S:=\{(x,y)\mid x^2+y^2=z^2\}$ 不是拓扑流形. 若它是拓扑流形, 则 $S$ 中存在一个含原点的连通开集 $W$, 与开子集 $U\subset\mathbb{R}^2$ 同胚. 此时 $W\setminus\{0\}$ 应与 $U\setminus\{q\}$ 同胚 (其中 $q$ 是原点的原像). 但是 $W\setminus\{0\}$ 不连通, 而 $U\setminus\{q\}$ 连通, 因而两者不可能同胚.
\end{enumerate}

\RGsolution{ex:klein-bottle-from-mobius-bands}{klein-bottle-from-mobius-bands}
容易证明, 从 Klein 瓶中删去一个 Möbius 带后, 所得空间又与一个 Möbius 带同胚 (见 Fig.~\ref{fig:klein-bottle-minus-mobius-band}).
\RGfigure[.62\textwidth]{klein-bottle-minus-mobius-band}{删去一个 Möbius 带的 Klein 瓶}

\RGsolution{ex:euler-characteristic-of-surfaces}{euler-characteristic-of-surfaces}
\begin{enumerate}[label=(\alph*),ref=\alph*,leftmargin=*]
\item\label{sol:euler-characteristic-invariance-under-adding-vertices} 若新顶点位于某个面上, 则必须将它与该面的三个顶点相连. 这样一来, 顶点数增加 1, 边数增加 3, 面数增加 2 (因为一个面被分成了三个面). 因此, $V-E+F$ 的变化量为 $1-3+2=0$. 若新顶点位于某条边上, 则必须将它与该边对面的两个顶点相连. 这样一来, 顶点数增加 1, 边数增加 3 (因为增添了两条边, 且一条边被分成了两条), 面数增加 2 (因为两个面被分成了四个面). 因而 $V-E+F$ 仍无变化.
\item 对由四面体确定的 $S^2$ 的三角剖分, 有 $V=4$, $E=6$ 且 $F=4$, 所以 $\chi(S^2)=4-6+4=2$.
\item 给边两两等同的正方形添上一条对角线, 即可将 $T^2$ 分解为三角形. 严格来说, 这不是一个三角剖分 (因为两个三角形的交集由全部三条边组成), 但可以通过增添顶点将它变成三角剖分; 由 (\ref{sol:euler-characteristic-invariance-under-adding-vertices}), 可以用这一分解计算 Euler 特征数. 计入边和顶点的等同关系后, 有 $V=1$, $E=3$, $F=2$, 因而 $\chi(T^2)=1-3+2=0$.
\item 与 $T^2$ 的情形相同.
\item 与 $T^2$ 的情形相同, 只不过此处等同后 $V=2$, 因而 $\chi(\mathbb{R}P^2)=2-3+2=1$.
\item 考虑 $M$ 和 $N$ 的三角剖分. 从两个曲面各删去一个三角形, 再把所删三角形的边对应等同, 即得到 $M\#N$ 及其一个三角剖分. 由于有三个顶点两两等同, 顶点总数减少 3, 边总数也减少 3. 面总数减少 2, 对应于所删去的两个面. 因此 $\chi(M\#N)=\chi(M)+\chi(N)-2$.
\end{enumerate}
\end{description}
